\documentclass{article}

\usepackage{amsmath,amssymb}
\usepackage{xurl}
\usepackage[hidelinks]{hyperref}
\setlength{\emergencystretch}{2em}

% Shared commands for notes in docs/paper-gaps/.

\DeclareUnicodeCharacter{2080}{\ifmmode{}_{0}\else\textsubscript{0}\fi}
\DeclareUnicodeCharacter{2081}{\ifmmode{}_{1}\else\textsubscript{1}\fi}
\DeclareUnicodeCharacter{2082}{\ifmmode{}_{2}\else\textsubscript{2}\fi}
\DeclareUnicodeCharacter{2099}{\ifmmode{}_{n}\else\textsubscript{n}\fi}

\newcommand{\Lean}{\textsc{Lean}}
\ExplSyntaxOn
\NewDocumentCommand{\leanid}{m}{
  \ifmmode
    \textsf{\detokenize{#1}}
  \else
    \group_begin:
    \urlstyle{sf}
    \tl_set:Nn \l_tmpa_tl {#1}
    \tl_replace_all:Nnn \l_tmpa_tl {\_} {_}
    \exp_args:NV \path \l_tmpa_tl
    \group_end:
  \fi
}
\ExplSyntaxOff
\providecommand{\tr}{\operatorname{tr}}
\newcommand{\tnleanIssueBase}{https://github.com/LionSR/TNLean/issues/}
\newcommand{\tnleanPrBase}{https://github.com/LionSR/TNLean/pull/}
\newcommand{\tnleanDocBase}{https://sirui-lu.com/TNLean/docs/}
\newcommand{\ghissue}[1]{\href{\tnleanIssueBase#1}{GitHub issue \##1}}
\newcommand{\ghpr}[1]{\href{\tnleanPrBase#1}{PR \##1}}
\newcommand{\doclink}[1]{\href{\tnleanDocBase#1.html}{\path{#1}}}

% Dirac notation and the identity operator, as in blueprint/src/macros/common.tex.
% \providecommand keeps note-local definitions authoritative.
\usepackage{dsfont}
\providecommand{\ket}[1]{|#1\rangle}
\providecommand{\bra}[1]{\langle #1|}
\providecommand{\braket}[2]{\langle #1 | #2 \rangle}
\providecommand{\ketbra}[2]{|#1\rangle\!\langle #2|}
\providecommand{\Id}{\mathds{1}}
\providecommand{\one}{\mathds{1}}
\providecommand{\C}{\mathbb{C}}
\providecommand{\MN}[1]{M_{#1}(\C)}
\providecommand{\MD}{M_D(\C)}

% External referencing for paper sources.  The build helper writes sanitized
% auxiliary files under build/paper-aux/ and excludes duplicated source labels.
\usepackage{xr}

% Import a generated paper auxiliary file with a caller-supplied label prefix.
% The candidate paths support builds from docs/paper-gaps/, the repository root,
% and build/paper-aux/ itself.
\newcommand{\importpaperaux}[2]{%
  \IfFileExists{../../build/paper-aux/#2.aux}{%
    \externaldocument[#1]{../../build/paper-aux/#2}%
  }{%
    \IfFileExists{build/paper-aux/#2.aux}{%
      \externaldocument[#1]{build/paper-aux/#2}%
    }{%
      \IfFileExists{#2.aux}{%
        \externaldocument[#1]{#2}%
      }{}%
    }%
  }%
}

% General external reference macros
\newcommand{\paperref}[2]{\ref{#1-#2}}
\newcommand{\papereqref}[2]{(\ref{#1-#2})}

% CPSV16 source (1606.00608 / MPDO-22-12-17-2.tex) external document and shortcuts
\importpaperaux{cpsv16-}{1606.00608/MPDO-22-12-17-2-tnlean}
\newcommand{\cpsvref}[1]{\paperref{cpsv16}{#1}}
\newcommand{\cpsveqref}[1]{\papereqref{cpsv16}{#1}}

% Verdict marker: \gapnote{<kind>}{<status>}. Kind is the policy
% classification (clarification, local-correction, scope-restriction,
% unfaithful, false-source, open-gap); status is open, wip (elimination
% underway), resolved, or historical. Machine-read by the site index;
% prints a verdict line.
\newcommand{\gapnote}[2]{\par\noindent\textbf{Verdict:} #1 (#2).\par}


\title{Shifted Trace Powers and the Exact MPU Transfer Spectrum}
\author{}
\date{2026-08-09}

\begin{document}
\maketitle
\gapnote{open-gap}{resolved}

\section{The source statement}

Let $U$ be a matrix-product unitary of physical dimension $d$ and bond
dimension $D$. Combine its two physical indices into
$a=(p,q)\in\{1,\ldots,d\}^2$, and define the normalized flattening by
\begin{align}
    \widetilde U^{(p,q)}
        &= d^{-1/2}U^{p,q}.
    \label{eq:cpsv17_transfer_trace_power_normalized_flattening}
\end{align}
The associated transfer map is the linear endomorphism
$T:M_D(\mathbb C)\to M_D(\mathbb C)$ defined by
\begin{align}
    T(X)
        &= \sum_{a\in\{1,\ldots,d\}^2}
           \widetilde U^a X(\widetilde U^a)^\dagger.
    \label{eq:cpsv17_transfer_trace_power_transfer_map}
\end{align}
This is \leanid{Kraus.transferMap} applied to
\leanid{MPOTensor.normalizedFlattening}.

Let $F_{k\ell}$ denote the matrix units in $M_D(\mathbb C)$, and let
$E=\widehat T$ be the column-stacked matrix representation of $T$. The
column-stacking convention is
\begin{align}
    \operatorname{vec}(X)_{(j,i)}
        &= X_{ij}.
    \label{eq:cpsv17_transfer_trace_power_column_stacking}
\end{align}
With this convention, $E$ is characterized by
\begin{align}
    E_{(j,i),(\ell,k)}
        &= (T(F_{k\ell}))_{ij},
    \label{eq:cpsv17_transfer_trace_power_transfer_matrix_entries}\\
    E\,\operatorname{vec}(X)
        &= \operatorname{vec}(T(X)).
    \label{eq:cpsv17_transfer_trace_power_transfer_matrix_action}
\end{align}
These identities give the defining convention of \leanid{transferMatrix} and
its application lemma \leanid{transferMatrix_mulVec_eq}.

The transfer matrix is indexed by pairs of bond indices and therefore has
matrix dimension
\begin{align}
    n
        &= D^2.
    \label{eq:cpsv17_transfer_trace_power_transfer_dimension}
\end{align}
For any $n\times n$ matrix $B$, define its characteristic polynomial by
\begin{align}
    \chi_B(X)
        &= \det(X\mathbf 1_n-B).
    \label{eq:cpsv17_transfer_trace_power_characteristic_polynomial}
\end{align}
In particular, $\chi_E$ is the characteristic polynomial of the transfer
matrix.

For every integer $N>1$, Proposition~III.1 of
Ref.~\cite{Cirac2017MPU} derives the shifted moment identity
\begin{align}
    \tr(E^N)
        &= 1,
        \qquad N>1.
    \label{eq:cpsv17_transfer_trace_power_shifted_moments}
\end{align}
The source concludes that $1$ is the only nonzero eigenvalue of $E$ and has
algebraic multiplicity one. Equivalently,
\begin{align}
    \chi_E(X)
        &= X^{n-1}(X-1).
    \label{eq:cpsv17_transfer_trace_power_source_characteristic_polynomial}
\end{align}
The cited Lemma~A.5 of Ref.~\cite{Cirac2016MPDO_arXiv} is naturally stated
using all positive moments, including the first. The MPU hypothesis supplies
Eq.~\ref{eq:cpsv17_transfer_trace_power_shifted_moments}, but it does not
supply $\tr(E)=1$.

\section{The shifted-moment argument}

Fix the transfer matrix $E$, its dimension $n$, and its characteristic
polynomial $\chi_E$. For every integer $m\geq1$,
Eq.~\ref{eq:cpsv17_transfer_trace_power_shifted_moments} gives
\begin{align}
    \tr((E^2)^m)
        &= \tr(E^{2m}),
    \label{eq:cpsv17_transfer_trace_power_square_moment_reduction}\\
    \tr(E^{2m})
        &= 1,
    \label{eq:cpsv17_transfer_trace_power_square_moments}\\
    \tr((E^3)^m)
        &= \tr(E^{3m}),
    \label{eq:cpsv17_transfer_trace_power_cube_moment_reduction}\\
    \tr(E^{3m})
        &= 1.
    \label{eq:cpsv17_transfer_trace_power_cube_moments}
\end{align}
The all-positive-moment characteristic-polynomial theorem applies separately
to $E^2$ and $E^3$, yielding
\begin{align}
    \chi_{E^2}(X)
        &= X^{n-1}(X-1),
    \label{eq:cpsv17_transfer_trace_power_square_characteristic_polynomial}\\
    \chi_{E^3}(X)
        &= X^{n-1}(X-1).
    \label{eq:cpsv17_transfer_trace_power_cube_characteristic_polynomial}
\end{align}

Let $\lambda\neq0$ be a spectral value of $E$. Spectral mapping and
Eqs.~\ref{eq:cpsv17_transfer_trace_power_square_characteristic_polynomial}
and~\ref{eq:cpsv17_transfer_trace_power_cube_characteristic_polynomial} give
\begin{align}
    \lambda^2
        &= 1,
    \label{eq:cpsv17_transfer_trace_power_square_spectral_constraint}\\
    \lambda^3
        &= 1,
    \label{eq:cpsv17_transfer_trace_power_cube_spectral_constraint}\\
    \lambda
        &= 1.
    \label{eq:cpsv17_transfer_trace_power_nonzero_spectral_value}
\end{align}
Conversely, the eigenvalue $1$ belongs to the spectrum of $E^2$ by
Eq.~\ref{eq:cpsv17_transfer_trace_power_square_characteristic_polynomial}.
Spectral mapping therefore supplies a nonzero spectral value of $E$. Hence
\begin{align}
    \sigma(E)\setminus\{0\}
        &= \{1\}.
    \label{eq:cpsv17_transfer_trace_power_nonzero_spectrum}
\end{align}
This set equality does not by itself determine algebraic multiplicity.

To control multiplicity, use the commuting factorization
\begin{align}
    E^2-\mathbf 1
        &= (E+\mathbf 1)(E-\mathbf 1).
    \label{eq:cpsv17_transfer_trace_power_square_factorization}
\end{align}
Every generalized $1$-eigenvector of $E$ is therefore a generalized
$1$-eigenvector of $E^2$. It follows that
\begin{align}
    \operatorname{mult}_1(\chi_E)
        &\leq \operatorname{mult}_1(\chi_{E^2}),
    \label{eq:cpsv17_transfer_trace_power_multiplicity_comparison}\\
    \operatorname{mult}_1(\chi_{E^2})
        &= 1.
    \label{eq:cpsv17_transfer_trace_power_square_multiplicity}
\end{align}
Because $1\in\sigma(E)$ by
Eq.~\ref{eq:cpsv17_transfer_trace_power_nonzero_spectrum}, its multiplicity is
positive. Equations~\ref{eq:cpsv17_transfer_trace_power_multiplicity_comparison}
and~\ref{eq:cpsv17_transfer_trace_power_square_multiplicity} therefore give
\begin{align}
    \operatorname{mult}_1(\chi_E)
        &= 1.
    \label{eq:cpsv17_transfer_trace_power_unit_multiplicity}
\end{align}
All other nonzero roots were excluded in
Eq.~\ref{eq:cpsv17_transfer_trace_power_nonzero_spectrum}. Since the monic
characteristic polynomial splits over $\mathbb C$, one obtains
\begin{align}
    \chi_E(X)
        &= X^{n-1}(X-1).
    \label{eq:cpsv17_transfer_trace_power_derived_characteristic_polynomial}
\end{align}
The argument uses no identity for $\tr(E)$ and requires no positivity,
normality, or diagonalizability assumption.

\section{Verdict and formalization status}

\textbf{Verdict: the shifted-moment multiplicity gap is closed.}
The set-spectrum result in
Eq.~\ref{eq:cpsv17_transfer_trace_power_nonzero_spectrum} remains available
as a deliberately narrower consequence. The exact source conclusion in
Eq.~\ref{eq:cpsv17_transfer_trace_power_derived_characteristic_polynomial} is
now formalized directly from moments with exponent greater than
one.\footnote{Formal declarations:
\leanid{Matrix.spectrum_diff_zero_eq_singleton_of_forall_trace_pow_eq_one_of_one_lt},
\leanid{Matrix.charpoly_rootMultiplicity_one_eq_one_of_forall_trace_pow_eq_one_of_one_lt},
and
\leanid{Matrix.charpoly_eq_X_pow_pred_mul_X_sub_one_of_forall_trace_pow_eq_one_of_one_lt}
in \doclink{TNLean/Algebra/ShiftedTracePowerSpectrum}.}
For the normalized flattening of an MPU with bond dimension $D$, the
specialization gives the exponent $D^2-1$.\footnote{Formal declaration:
\leanid{MPOTensor.IsMPU.normalizedFlattening_charpoly} in
\doclink{TNLean/MPS/MPU/TransferMatrix}.}

\bibliographystyle{alpha}
\bibliography{references}

\end{document}
